\\
Fix $k\in\mathbb{N}$. For an integer $n\ge 1$, let
\[
 j(n) := \min\{m\in\mathbb{N}: \text{every interval of $m$ consecutive integers contains an integer coprime to $n$}\}.
\]
(Here ``interval of $m$ consecutive integers'' means a set of the form $\{t+1,t+2,\dots,t+m\}$ for some $t\in\mathbb{Z}$.  ``Coprime'' means $\gcd(\cdot,\cdot)=1$.)
Define
\[
 h(k) := \max\{ j(n): n\ge 1\text{ and $n$ has at most $k$ prime factors}\}.
\]
\textbf{Ambiguity.} ``At most $k$ prime factors'' can mean either $\Omega(n)\le k$ (counted with multiplicity) or $\omega(n)\le k$ (distinct prime divisors).  In either convention, the value of $j(n)$ depends only on the distinct prime divisors of $n$ (Lemma~\ref{lem:970-rad}), so $h(k)$ is unchanged by replacing $\Omega$ with $\omega$.

The question asks for the order of magnitude of $h(k)$ as $k\to\infty$; in particular whether there exists an absolute constant $C$ such that $h(k)\le C k^2$ for all sufficiently large $k$.
